% =====================================================================
%  ANEXO B: GUÍA DE ANOTACIÓN
% =====================================================================
\anexo{B}{Guía de Anotación}

La etiqueta describe lo que cuenta el relato; no confirma hechos ni señala responsables. Los anotadores no deben suponer el género, la orientación, el consentimiento ni la intención de nadie a partir de estereotipos.

\begin{table}[H]
\begin{threeparttable}
\caption{Etiquetas y criterios de anotación}
\label{tab:anotacion}
\renewcommand{\arraystretch}{1.3}
\begin{tabular}{@{}L{3.3cm} L{7.4cm} L{4.6cm}@{}}
\toprule
\textbf{Etiqueta} & \textbf{Criterio} & \textbf{Tratamiento} \\
\midrule
Con acoso sexual & El relato describe una conducta sexual no deseada, una presión o un condicionamiento sexual, con contexto suficiente & Clase positiva (1) \\
Sin acoso sexual & El relato tiene contexto suficiente y no cumple los criterios anteriores & Clase negativa (0) \\
Información insuficiente & No se puede saber qué ocurrió, en qué contexto o si se trata de una cita & Se revisa con un especialista; si persiste la duda, se excluye \\
\bottomrule
\end{tabular}
\begin{tablenotes}[flushleft]
\small
\item \textit{Nota.} Elaboración propia.
\end{tablenotes}
\end{threeparttable}
\end{table}

\noindent\textbf{Reglas de decisión:}
\begin{itemize}
	\item Distinguir entre lo que el estudiante vivió y lo que cita o cuenta de otros.
	\item Revisar las negaciones y las menciones a hechos que no ocurrieron.
	\item Tomar en cuenta una relación de dependencia (docente, jefe de práctica, tutor) solo si aparece en el relato.
	\item No exigir que la conducta se repita ni que el agresor tenga más jerarquía.
	\item No etiquetar como acoso una conversación educativa sobre sexualidad solo por las palabras que usa.
	\item Ante la duda, marcar \textit{información insuficiente} en lugar de forzar una etiqueta.
\end{itemize}
